\chapter{Introduction}\label{ch:introduction}

\section{Background and Motivation}\label{sec:intro_background}
Large Language Models built on the transformer
architecture~\cite{vaswani2017attention} have, in the span of a handful
of years, moved from a research curiosity to the primary workload of
most new supercomputing deployments. The trajectory from GPT-3's
$175$-billion parameters~\cite{brown2020gpt3} through PaLM's
$540$-billion~\cite{chowdhery2022palm} and the open-weight LLaMA
series~\cite{touvron2023llama,touvron2023llama2} is less a qualitative
jump at each step than a straightforward consequence of the scaling
laws~\cite{kaplan2020scaling,hoffmann2022chinchilla}: for any given
compute budget the loss-minimising model is larger and trained on
more data than the previous-generation wisdom suggested. The
engineering consequence is that every new model generation
simultaneously tightens the memory budget of a single GPU, widens
the communication demand across a distributed cluster, and
lengthens the context windows on which the attention operator is
evaluated~\cite{rajbhandari2019zero,shoeybi2019megatron}. A software
stack that was adequate for the previous model generation is
rarely adequate for the next one without retuning at every layer.

\subsection{A Concrete Memory Accounting}\label{sec:intro_memory_example}
A single back-of-the-envelope estimate makes the stakes tangible. For a
$7$-billion-parameter transformer trained with the Adam optimiser in FP16
mixed precision, the per-parameter memory footprint is
$2$\,B (parameter) $+ 2$\,B (gradient) $+ 4 \cdot 2$\,B (FP32 master copy
and the two Adam moments) $= 14$\,B, or roughly $98$\,GB just to hold the
model state. A single A100-$40$\,GB cannot hold this state---let alone
the activations, which at sequence length $N{=}2048$ and batch $B{=}4$
add a further $\sim\!20$\,GB per layer. A $70$-billion-parameter model
multiplies the state budget by ten and the activation budget by three,
because $d_{\text{model}}$ and the number of layers both grow. Memory
is therefore the first-order constraint on \emph{what can be trained at
all}, and it is the constraint that drives every sharding decision in
this thesis.

\subsection{The Four Challenges That Motivate This Thesis}\label{sec:intro_four_challenges}

Modern LLMs require resources that far exceed the capability of a single
GPU. Training models such as PaLM or LLaMA-3 necessitates distributed
computing across hundreds or thousands of accelerators. Four interacting
challenges arise as soon as one moves beyond a single GPU:
\begin{itemize}
    \item \textbf{Memory management.} Model parameters, gradients, and
          optimiser states for a $7$-billion parameter LLM already exceed
          the $40$\,GB HBM of a single A100. Adam~\cite{rasley2020deepspeed}
          requires two additional optimiser tensors per parameter (first
          and second moments), tripling the parameter memory; mixed
          precision training~\cite{micikevicius2018amp} requires an
          additional FP32 master copy.
    \item \textbf{Communication overhead.} Gradient synchronisation over a
          data-parallel group scales with the model size, not the batch
          size. Without careful overlap of communication and compute, the
          accelerators spend a majority of their time waiting on NVLink
          or InfiniBand.
    \item \textbf{Hardware heterogeneity.} Production stacks must support
          diverse hardware (A100, H100, Intel Xe, AMD CDNA) with different
          memory hierarchies, Tensor-Core instructions, and interconnects.
          The same training job can exhibit markedly different performance
          profiles on different SKUs.
    \item \textbf{Attention efficiency.} The $\mathcal{O}(N^2)$ complexity
          of self-attention dominates the FLOPs of long-context models.
          Simple masking of an element-wise sparse pattern is insufficient,
          because dense GEMMs are still issued over blocks that are
          subsequently zeroed out.
\end{itemize}

The carbon and energy cost of training~\cite{patterson2021carbon} further
amplifies the importance of efficiency: any $10\%$ improvement in compute
utilisation translates directly into a $10\%$ reduction in electricity
consumption and in embodied carbon. Optimisations that simultaneously
reduce wall-clock time and energy per epoch are therefore valuable beyond
their direct performance benefit.

\subsection{A Stack View of the Problem}\label{sec:intro_stack}
Throughout this thesis we treat LLM efficiency not as a single
algorithmic question but as a \emph{stack} of co-dependent layers, each
of which must be tuned against the layer above and the layer below:
\begin{enumerate}
    \item \textbf{Distributed schedule} (Chapter~\ref{ch:distributed}).
          Decides which parts of the model, gradients, and optimiser
          state live on which GPU, and which collective operations keep
          them synchronised. Sets the memory envelope for every kernel
          call below it.
    \item \textbf{Kernel layer} (Chapter~\ref{ch:cutlass}). Provides the
          dense GEMMs, grouped GEMMs, and epilogue-fused primitives on
          which every transformer layer is built. Its Tensor-Core
          utilisation is inherited at every scale of parallelism.
    \item \textbf{Attention algorithm} (Chapter~\ref{ch:swfa}).
          Determines the asymptotic complexity of the dominant operator
          at long context. A kernel that ignores the sparsity pattern
          exposed by the model architecture wastes the savings the
          upper layers have worked to create.
    \item \textbf{Serving runtime} (Chapter~\ref{ch:sglang}). Turns
          algorithmic savings into end-user-visible throughput by
          scheduling kernel launches, managing the KV cache, and
          batching requests continuously.
    \item \textbf{Profiling and auto-tuning}
          (Chapter~\ref{ch:triton}). Closes the loop: it measures
          whether the four layers above are actually delivering their
          expected performance, and it selects per-shape kernel
          parameters that would otherwise have to be hand-tuned.
\end{enumerate}
This view has two consequences that inform every chapter. First, a
point optimisation at one layer is only as useful as the layer above
allows it to be: a kernel speedup disappears if the distributed
schedule serialises it behind a slow collective. Second, a measurement
at one layer is only as trustworthy as the cross-check at the layer
below: a higher-level throughput number can hide a kernel regression
that a lower-level profiler would have caught. The profiling pipeline
of Chapter~\ref{ch:triton} is therefore treated as a first-class
deliverable, not as a side-tool.

\section{Problem Statement}\label{sec:intro_problem}
The primary problems addressed in this thesis are as follows.
\begin{enumerate}
    \item \textbf{Scalable distributed training.} Existing comparisons of
          DP, DDP, FSDP, and DeepSpeed ZeRO-3 are typically published by
          the authors of each system, each under a different workload and
          a different hardware configuration. A controlled single-node
          study that fixes the model, the dataset, the random seed, and
          the hardware, and varies only the parallelisation strategy, is
          missing from the literature.
    \item \textbf{High-performance matrix multiplication.} Vendor cuBLAS
          kernels are highly tuned for large square GEMMs but are
          suboptimal for skinny, batched, and grouped shapes---precisely
          the shapes that dominate transformer attention and
          Mixture-of-Experts layers. A principled mapping of LLM GEMMs
          onto the CUTLASS/CuTe abstraction is required.
    \item \textbf{Attention optimisation.} Sliding-window
          attention~\cite{beltagy2020longformer} has become ubiquitous in
          long-context models, yet most open-source implementations apply
          the window as a post-hoc element-wise mask. A kernel that
          exploits the window structure \emph{before} computing scores
          would yield near-linear scaling in the window size.
    \item \textbf{Production-grade serving kernels.} LLM inference is
          qualitatively different from training: requests arrive
          asynchronously, prompts share large prefixes, and the decode
          phase is memory-bound. Kernel optimisations for the SGLang
          runtime---paged KV cache, fused decode attention, fused
          sampling---are needed to turn algorithmic savings into wall-clock
          serving throughput.
    \item \textbf{Kernel profiling and tuning.} Kernel-level performance
          decisions (block sizes, pipeline stages, warp counts) need to be
          revisited for every new GPU generation and every new input
          shape. A uniform profiling pipeline that crosses Triton,
          CUTLASS, and Torch-inductor is required to keep these decisions
          current.
\end{enumerate}

\section{Objectives}\label{sec:intro_objectives}
\begin{enumerate}
    \item Design and evaluate a modular high-performance deep-learning
          software stack that supports efficient LLM training and
          inference across diverse hardware, with clean interfaces between
          the distributed scheduler, the kernel layer, the attention
          algorithm, and the serving runtime.
    \item Implement and benchmark four distributed training
          strategies---DP, DDP, FSDP, and DeepSpeed ZeRO-3---on an 8-GPU
          server, and report GPU utilisation, per-GPU memory footprint,
          training-loss convergence, power consumption, energy per epoch,
          and scaling efficiency on a single unified workload.
    \item Study CUTLASS/CuTe-based GEMM pipelines and map the principal
          LLM layers (QKV, MLP, attention score, attention output, grouped
          MoE) onto the Tensor-Core hierarchy, reporting Tensor-Core
          utilisation and end-to-end speedups against cuBLAS.
    \item Develop a block-sparse Sliding Window Flash Attention (SWFA)
          algorithm with a rigorous accuracy guarantee, a closed-form
          optimal-window derivation, and a numerical validation against a
          vanilla attention reference over nineteen configurations.
    \item Engineer GPU kernels for the SGLang serving runtime---paged KV
          cache, decode attention, prefill attention, and fused
          sampling---and measure their cumulative effect on LLaMA-7B
          decode throughput.
    \item Build a Triton-based profiling workflow for kernel-level
          analysis that combines CUDA-event timing, roofline analysis,
          and Nsight Compute cross-checks, with an auto-tuner that caches
          configurations keyed by input shape and data type.
\end{enumerate}

\section{Contributions}\label{sec:intro_contributions}
The main contributions of this thesis are as follows.
\begin{itemize}
    \item \textbf{Controlled distributed-training study.} A systematic,
          controlled comparison of four distributed training strategies
          on an 8-GPU A100 server, with detailed metrics for GPU
          utilisation, memory, loss convergence, power consumption, and
          scaling efficiency (Chapter~\ref{ch:distributed}). The
          measurements show that ZeRO-3 delivers $22\%$ lower per-epoch
          energy than DP at the same final loss.
    \item \textbf{CUTLASS/CuTe mapping.} A mapping of the dense LLM
          computations (QKV, attention score/output, MLP, MoE) onto
          CUTLASS/CuTe GEMM pipelines, covering tile sizes, software
          pipelining, Tensor-Core utilisation, and epilogue fusion
          (Chapter~\ref{ch:cutlass}).
    \item \textbf{Sliding Window Flash Attention.} A novel SWFA algorithm
          with an efficient block-classification step and a vectorised
          mask generator that achieves up to $65.6\%$ reduction in
          computed blocks while preserving numerical accuracy within
          $10^{-4}$ tolerance (Chapter~\ref{ch:swfa}). A closed-form
          expression for the optimal window size under an exponential
          attention-decay model is also derived; strict convexity of the
          associated loss is proved.
    \item \textbf{SGLang kernel stack.} GPU kernel optimisations inside
          the SGLang runtime for long-context LLM serving, including
          paged KV cache, fused decode attention, SWFA prefill, and fused
          top-$k$/top-$p$ sampling. A layered evaluation shows a
          cumulative $2.84\times$ decode throughput speedup on LLaMA-7B
          (Chapter~\ref{ch:sglang}).
    \item \textbf{Triton profiling pipeline.} A profiling and auto-tuning
          pipeline for kernel-level roofline analysis, with shape-indexed
          configuration caching and a cross-validator that reports
          correctness and throughput against CUTLASS and Torch-inductor
          (Chapter~\ref{ch:triton}).
    \item \textbf{Reproducibility.} Algorithms, mask generators, and the
          block classifier are given as pseudocode in the main text and
          as reference PyTorch implementations in Appendix~A; full proofs
          of the convexity and optimal-window results, together with the
          block-count and communication-cost lemmas, are collected in
          Appendix~B.
\end{itemize}

\section{Scope of the Work}\label{sec:intro_scope}
The experimental evaluation in this thesis focuses on NVIDIA A100 GPUs in
an 8-GPU NVLink-connected server, with cross-checks on CPU where relevant
(notably for the SWFA numerical accuracy tests). Distributed training is
evaluated on ResNet-50/CIFAR-10 as a controlled proxy and on a GPT-2
Medium-class model ($355$M parameters) for end-to-end LLM throughput
measurements. The SWFA kernel is implemented in PyTorch with CPU-optimised
paths and is prototyped for a later CUTLASS/Triton port; correctness is
established on the reference implementation.

The work deliberately does not cover: (i) multi-node scaling, which is
left as future work in Chapter~\ref{ch:conclusion}; (ii) inference
quantisation below $8$ bits, which is orthogonal to the kernel-level
optimisations studied here; and (iii) reinforcement-learning-from-human-%
feedback (RLHF) or alignment workloads, which exercise different
bottlenecks. Where relevant, pointers to the relevant literature are
given so that the findings of this thesis can be placed in the broader
context.

\section{Methodology Overview}\label{sec:intro_methodology}
Each technical chapter (Chapters~\ref{ch:distributed}--\ref{ch:triton})
follows a uniform internal structure that makes it possible to read any
chapter in isolation while still understanding how it fits into the
overall stack:
\begin{enumerate}
    \item \textbf{Background.} A concise restatement of the relevant
          concepts, calibrated to the specific layer of the stack.
    \item \textbf{Methodology.} An explicit \emph{Methodology} section
          describing the experimental setup, the workload, the metrics,
          and the hardware/software versions. The same hardware
          description is used across chapters so that numbers from
          different chapters can be compared directly.
    \item \textbf{Implementation.} Either pseudocode or reference
          PyTorch/Triton snippets that are sufficient for an independent
          reimplementation.
    \item \textbf{Results.} Quantitative results on the workload of the
          chapter, always accompanied by tables and figures whose source
          data is reproducible from the reference implementation.
    \item \textbf{Discussion.} A brief discussion of what the numbers
          mean and how they compose with the rest of the stack.
\end{enumerate}
This structure keeps the experimental setup alongside the measurements it
produces and avoids a single monolithic methodology chapter that would
otherwise span five unrelated systems.

\section{Timeline}\label{sec:intro_timeline}
The work reported in this thesis was carried out in four phases:
\begin{enumerate}
    \item \textbf{Phase 1 --- Distributed training (months 1--3).}
          Baseline DP/DDP/FSDP/ZeRO-3 implementations, setup of the
          8-GPU A100 server, NCCL/DeepSpeed tuning, and the controlled
          comparison of Chapter~\ref{ch:distributed}.
    \item \textbf{Phase 2 --- Kernel layer (months 3--6).} CUTLASS/CuTe
          study of Chapter~\ref{ch:cutlass}, including the LLM-layer
          mapping and epilogue fusion micro-benchmarks.
    \item \textbf{Phase 3 --- Attention algorithm (months 6--9).}
          Design and implementation of SWFA
          (Chapter~\ref{ch:swfa}), numerical validation against a
          vanilla reference, and the optimal-window derivation.
    \item \textbf{Phase 4 --- Serving and profiling (months 9--12).}
          SGLang kernel engineering
          (Chapter~\ref{ch:sglang}) and the Triton profiling pipeline
          (Chapter~\ref{ch:triton}), which together close the loop
          between algorithmic design and on-GPU performance.
\end{enumerate}

\section{Thesis Organisation}\label{sec:intro_organisation}
Figure~\ref{fig:intro_stack} summarises the logical dependency between
the chapters. Each technical chapter both \emph{consumes} the
abstraction below it and \emph{produces} an abstraction for the chapter
above: for example, Chapter~\ref{ch:cutlass} delivers a GEMM primitive
that Chapter~\ref{ch:swfa} uses inside its block-wise attention kernel,
and Chapter~\ref{ch:sglang} in turn plugs that attention kernel into a
full serving runtime.

\begin{figure}[H]
    \centering
    \begin{tikzpicture}[
        node distance=3mm and 8mm,
        lyr/.style={draw, rounded corners, align=center,
                    minimum width=75mm, minimum height=8mm,
                    font=\small\bfseries},
        arr/.style={-latex, thick}
    ]
    \node[lyr, fill=blue!8]  (sched)   {Ch.~3 --- Distributed schedule
                                        (DP / DDP / FSDP / ZeRO-3)};
    \node[lyr, fill=blue!10, below=of sched] (kern)    {Ch.~4 --- CUTLASS /
                                        CuTe kernel primitives};
    \node[lyr, fill=blue!12, below=of kern]  (attn)    {Ch.~5 --- SWFA
                                        attention algorithm};
    \node[lyr, fill=blue!14, below=of attn]  (serve)   {Ch.~6 --- SGLang
                                        serving runtime};
    \node[lyr, fill=blue!16, below=of serve] (prof)    {Ch.~7 --- Triton
                                        profiling \& auto-tune};
    \draw[arr] (sched) -- (kern);
    \draw[arr] (kern)  -- (attn);
    \draw[arr] (attn)  -- (serve);
    \draw[arr] (serve) -- (prof);
    \draw[arr, dashed] (prof.east) .. controls +(2.5,0) and +(2.5,0)
                      .. node[midway, right, font=\footnotesize]
                             {measured feedback} (sched.east);
    \end{tikzpicture}
    \caption{The proposed LLM software stack. Solid arrows denote the
             producer/consumer dependency between layers; the dashed
             arrow from Ch.~7 back to Ch.~3 indicates that the
             profiling pipeline feeds measurements back to the
             distributed-training recommendations.}
    \label{fig:intro_stack}
\end{figure}

\noindent The remainder of the thesis is organised as follows.
\begin{itemize}
    \item \textbf{Chapter~\ref{ch:litreview} --- Literature Review.}
          Surveys prior work on distributed training, GEMM libraries,
          efficient attention, LLM serving systems, and GPU kernel
          profiling. Positions this thesis against the open gaps in
          that literature and introduces the vocabulary used
          throughout.
    \item \textbf{Chapter~\ref{ch:distributed} --- Distributed
          Computing and Training.} Presents the distributed-training
          methodology and controlled comparison of DP, DDP, FSDP, and
          DeepSpeed ZeRO-3 on an 8-GPU A100 node. Establishes which
          parallelisation strategy maximises GPU utilisation and
          minimises energy per epoch on a fixed workload, and produces
          the per-layer memory budget that the rest of the stack must
          respect.
    \item \textbf{Chapter~\ref{ch:cutlass} --- CUTLASS and CuTe GEMM
          Operations.} Describes GEMM operations built on CUTLASS and
          the CuTe layout algebra, and maps the dense LLM
          computations (QKV, attention score/output, MLP, grouped-MoE)
          onto the Tensor-Core hierarchy. The epilogue fusions
          introduced here reappear in Chapters~\ref{ch:swfa}
          and~\ref{ch:sglang} as the writeback path of the
          attention-output GEMM and the logit-projection GEMM,
          respectively.
    \item \textbf{Chapter~\ref{ch:swfa} --- Sliding Window Flash
          Attention.} Introduces the SWFA algorithm, including
          block classification, vectorised mask generation, forward
          and backward passes, and a closed-form optimal-window
          derivation with a proof of strict convexity. SWFA consumes
          the GEMM primitives of Chapter~\ref{ch:cutlass} and is later
          plugged into the prefill kernel of
          Chapter~\ref{ch:sglang}.
    \item \textbf{Chapter~\ref{ch:sglang} --- GPU Kernels in the
          SGLang Runtime.} Describes the serving-side component of
          the stack: paged KV cache, fused decode attention, SWFA
          prefill, and fused top-$k$/top-$p$ sampling. A layered
          evaluation decomposes the cumulative decode-throughput
          speedup into per-optimisation contributions.
    \item \textbf{Chapter~\ref{ch:triton} --- Triton Kernel Profiling
          and Auto-tuning.} Covers the Triton-based profiling and
          auto-tuning pipeline that cross-validates the CUTLASS
          numbers of Chapter~\ref{ch:cutlass}, the SWFA numbers of
          Chapter~\ref{ch:swfa}, and the serving numbers of
          Chapter~\ref{ch:sglang}. The shape-indexed configuration
          cache introduced here is the mechanism by which the
          kernel-level decisions remain valid as shapes and
          hardware generations change.
    \item \textbf{Chapter~\ref{ch:conclusion} --- Conclusions and
          Future Work.} Summarises the contributions, distils the
          lessons that composed across chapters, and outlines
          short-, medium-, and long-term research directions.
\end{itemize}

Algorithmic pseudocode, reference PyTorch listings, and
reproduction-ready hardware/software details are collected in
Appendix~A; full proofs of the theoretical results stated in the main
text are collected in Appendix~B. Each result in the main text is
cross-linked to its proof, and each algorithm is cross-linked to its
reference implementation, so that any claim in the thesis can be
traced to its evidence in at most two clicks.
